%! ~~~ Packages Setup ~~~ 
\documentclass[]{../../../../tex/classes/styledArticle}
\usepackage{../../../../tex/packages/hebrewSupport}
\usepackage{../../../../tex/packages/mathShortcuts}
\usepackage{../../../../tex/packages/theoremsSupport}
\usepackage{../../../../tex/packages/enfancysections}

%! ~~~ Document ~~~

\author{שחר פרץ}
\title{\textit{לינארית 2א $\sim$ סיכום שביעי $\sim$ פולינומים וכו'}}
\date{28 לאפריל 2025}
\begin{document}
	\maketitle
	הערות: אנחנו סימנו אידיאל ב־$aR$ ובקורס מסמנים $Ra$, באופן כללי אפשר לדבר על אידיאל שמאלי ואידיאל ימני. תזכורת: $I \subseteq R$ היא אידיאל אם היא סגורה לחיבור ומקיימת את תכונת הבליעה. בתחום ראשי כל אידיאל הוא אידיאל ראשי. 
	
	\theo{ב־$\{0\} \neq R$ תחום ראשי אז כל אי פריק הוא ראשוני. }\begin{proof}
		יהי $o$ אי פריק (א"פ). יהיו $a, b \in \R$ כך ש־$p \mid ab$. תיקון [משבוע שעבר]: במקום $I = Ra + Rb$, נשתמש ב־$I = Ra + Rp$. 
		בכלל ש־$R$ תחום ראשי, קיים $c \in R$ כך ש־$I = Rc$, ו־$a, p \in I$ כלומר $c \mid a \land c \mid p$. $p$ א"פ ולכן $c \sim p$ או $c$ הפיך. 
		\begin{itemize}
			\item $c$ הפיך $\impliedby$ $ \in I$ $\impliedby$ $R = R \cdot 1 \in I \subseteq R$ $\impliedby$ $I = R$. קיימים $r, s \in R$ כך ש־$ra + sp = 1$. נכפיל ב־$b$ ונקבל $rab + spb = b$ וסה"כ $p \mid b$. 
			\item אם $c \sim p$, אז $p \mid c \land c \mid a$ ולכן $p \mid a$. 
		\end{itemize}
	\end{proof}
	
	\cola{אם $R$ תחום שלמות ראשי אזי יש פריקות יחידה למכפלה של אי פריקים עד כדי חברות. }
	
	\theo{יהיו $a, b \in R$, אז $a, b$ ייקראו זרים אם $\forall c \in R\co c \mid a \land c \mid b \implies c \in R^x$}
	\defi{יהי $g \in R$ כך ש־:
	\begin{enumerate}
		\item $g \mid a \land g \mi b$
		\item $\forall \ml \in R \co \ml \mid a \land \ml \mid b$
		\item $\ml \mid g$
	\end{enumerate}
	אז $g$ כנ"ל הוא הגורם המשותף המקסימלי של $a, b$, הוא $\gcd(a, b)$}
	
	\theo{יהי $R$ תחום שלמות ויהיו $a, b \in R$. נניח שקיימים $r, s \in R$ כך ש־$g = ra + sb$ אשר מחלק את $a, b$. אז: 
	\begin{itemize}
		\item $\gcd(a, b) = g$
		\item ה־$\gcd$ מוגדר ביחידות עד לכדי חברות. 
		\item בתחום ראשי, לכל $a, b$ קיים $g$ כנ"ל. 
	\end{itemize}}
(הערה: רק 3 באמת דורש תחום ראשי)
	\begin{proof}
		\begin{itemize}
			\item יהי $\ml \mid a, b$ אז $\ml \mid ra, sb$ וסה"כ $\ml \mid g$. 
			\item מ־1 (בערך) אם $g, g'$ מקיימים את היותם $\gcd$ אז $g' \mid g \land g' \mid g$ ולכן $g \sim g'$. 
			\item נסמן $I = Ra + Rb$. אז $I = Rg$, וקיימים $r, s \in R$ כך ש־$ra + sb = g$ ולכן $a, b \in I$ אז $g \mid a, b$ וסיימנו מ־$1$. 
		\end{itemize}
	\end{proof}
	\cola{בתחום ראשי, אם $a, b$ זרים אז $\exists r, s \in \R\co ra + sb = 1$ (אלגוריתם אוקלידס המורחב). }
	\theo{$\F[x]$ תחום ראשי. }
	\begin{proof}
		יהי $I \subseteq \F[x]$ אידיאל. אם $I = \{0\}$, הוא ראשי. אחרת, $I \neq \{0\}$, ואז: יהי $0 \neq p \in I$ פולינום מדרגה מינימלית, ויהי $f \in I$. אז קיימים $q, r \in \F[x]$ כך ש־$f = qp + r$. ידוע $\deg r < \deg p$. בגלל ש־$f \in I \land p \in I$ אז $f - qp \in I$. אם $r$ אינו $0$, קיבלנו סתירה למינימליות הדרגה של $p$. 
	\end{proof}
	הוכחה זהה עובדת בשביל להראות ש־$\Z$ תחום ראשי, אך עם דרגה במקום ערך מוחלט. 
	
	\defi{תחום שלמות נקרא \textit{אוקלידי} אם קיימת $N \co R \setminus \{0\} \to \Z$ כך ש־$\forall a, b \in R\setminus \{0\} \, \exists u, r \in R \co a = ub + r$ כאשר $r = 0$ או $N(b) > N(r)$}
	ברגע שיש לנו את ההגדרה של תחום אוקלידי, $N$ הפונקציה שתשתמש אותנו בשביל להראות את ההוכחה שכל תחום אוקלידי הוא תחום ראשי (בדומה לערך מוחלט או $\deg$ בהוכחות קודמות). ההפך אינו בהכרח נכון. 
	
	\defi{נורמה היא פונקציה $N \co R \to \Z$ כך שהיא סאב־אדטיבית $N(a + b) \le N(a) + N(b)$, כיפלית ו־$N(1) =1$. }
	
	\textbf{דוגמה }\textit{(חוג השלמות של גאוס)}\textbf{. }
	\[ R = \Z[i] = \{a + bi \mid a, b \in \Z\} \]
	הנורמה: 
	\[ N \co \R\to \Z_{\ge 0}, \ N(a + bi) = a^2 + b^2 = |a + bi|^2 \]
	בדומה להוכחה לפיה הערך המוחלט של מורכב הוא כפלי, ניתן להראות ש־$N$ כפלית. מי הם ההפיכים ב־$\Z[i]$? מי שמקיים $\ag \bg = 1$, כלומר: 
	\[ N(\ag)N(\bg) = 1 = N(1), \ \ag = a + bi, \ a^2 + b^2 = 1 \implies a + bi = \pm 1, \pm i \]
	
	\theo{יהי $p \in \Z$ ראשוני. התנאים הבאים שקולים: 
	\begin{itemize}
		\item $p$ פריק ב־$\Z[i]$
		\item $p = m^2 + n^2$ עבור $n, m \in \Z$
		\item $p = 2$ או $p \equiv 1 \pmod{4} $
		\item קיימים $r, s \in R$ כך ש־$ra + sb = 1$
	\end{itemize}}

	שימו לב ש־$\Z$ בתוך $\Z[i]$ לא סגורים לבליעה. 
	
	
	
	\ndoc
\end{document}
